\scanpage{20}
Ainsi les classes de $\Pi_I^{\unsure{\circ|0}}$ sont les classes de $\Pi_J^\circ$ distinctes de $\deleted{\Pi}U_{J,i}$, et de plus les \unsure{clas.|deux.} parmi
\[
U_{J,i}\cap\Sigma_{\unsure{\deleted{J}i|j|i}}^+,\ U_{J,i}\cap\Sigma_i^-
\]
qui sont non vides. Or les deux classes en question sont non vides, car \deleted{$\unsure{\Sigma_i|\Sigma_j}$} $U_{J,i}$ est un voisinage de $\deleted{\Sigma}E_i$, et
\[
\overline{\Sigma_i^+}\cap\deleted{\Sigma}E_i\ne\varnothing,\quad
\overline{\deleted{\unsure{E}}\Sigma_i^-}\cap\deleted{\Sigma}E_i\ne\varnothing.
\]
On \unsure{sait|voit} donc que le nb des classes est en \deleted{\unsure{tout|est}} $(n-1)+2=n+1=\operatorname{card} I+1$ (ce qui achève \unsure{aussi|ainsi} la dém. par récurrence sur le nb des classes).

\uline{Lemme} Toute classe de $\Pi_I^\circ$ est de la forme $X\setminus\bigcup_{\alpha\in\deleted{I}K}(\overbrace{\Sigma_\alpha^+\cup E_\alpha}^{\inserted{\widehat{\Sigma}_\alpha^+}})$ où $K$ est une partie \inserted{\deleted{non vide}} de $I$, et
% slanted
\reversemarginpar\margin{\shortstack[l]{$K$ est $\ne\varnothing$\\si $I\ne\varnothing$}}\normalmarginpar
$\forall\alpha\in K$, \deleted{on a} $\Sigma_\alpha^+\deleted{\unsure{\Rightarrow|\to}}\in\pi_0(X\setminus E_\alpha)$, tels que $\forall\alpha,\beta\in K$, $\alpha\ne\beta$, on ait $\widehat{\Sigma}_\alpha^+\cap\widehat{\Sigma}_\beta^+=\varnothing$ $(\Longleftrightarrow\Sigma_\alpha^+\cap\Sigma_\beta^+=\varnothing)$.
% archive sheet number: 19
